\chapter{ACL Reconstruction}

\section*{Introduction \& Clinical Indications}
Anterior cruciate ligament (ACL) reconstruction is a prevalent orthopedic surgical procedure aimed at restoring knee stability by replacing a torn anterior cruciate ligament with a graft. The ACL is a vital intra-articular ligament that plays a crucial role in stabilizing the knee joint by preventing anterior translation of the tibia relative to the femur and resisting rotational forces. Injuries to the ACL are commonly associated with high-impact sports activities, particularly those involving sudden stops, pivots, or changes in direction, leading to symptoms such as knee instability, episodes of "giving way," and increased risk of secondary injuries to the menisci and articular cartilage. The primary goal of ACL reconstruction is to re-establish the biomechanical integrity of the knee, allowing patients to return to their pre-injury activity levels while minimizing the risk of long-term degenerative changes associated with chronic instability.

\begin{itemize}
  \item ACL repair with graft.
  \item Knee ligament reconstruction.
  \item Hamstring autograft ACL reconstruction.
  \item Bone-patellar tendon-bone (BTB) ACL reconstruction.
  \item Quadriceps tendon autograft ACL reconstruction.
  \item Allograft ACL reconstruction.
\end{itemize}

The fundamental principle of ACL reconstruction is to replace the compromised native ligament with a robust tendon graft that is biologically incorporated into the bone and biomechanically replicates the function of the original ACL. This involves creating bone tunnels in the femur and tibia at the anatomical attachment sites of the native ACL, through which the harvested graft is passed. The graft is then secured under appropriate tension using various fixation devices, such as interference screws, cortical buttons, or staples, to achieve immediate stability. The rationale behind this approach is to restore both anterior-posterior and rotational stability to the knee joint, thereby preventing recurrent episodes of "giving way" that can lead to further meniscal tears, articular cartilage damage, and the premature onset of osteoarthritis.

The choice of graft material, whether autograft (from the patient's own body, such as patellar tendon, hamstring tendons, or quadriceps tendon) or allograft (from a cadaveric donor), is a critical decision influenced by patient factors, surgeon preference, and activity level. The technical goals include achieving isometric graft placement, ensuring adequate graft tension, and promoting biological healing of the graft within the bone tunnels. By meticulously recreating the anatomical course and tension of the ACL, the reconstructed ligament provides the necessary mechanical restraint, allowing the patient to safely return to high-demand activities and protecting the long-term health of the knee joint.

The history of ACL reconstruction has evolved significantly over the past century. Early attempts at ACL repair in the late 19\textsuperscript{th} and early 20\textsuperscript{th} centuries involved direct repair of the torn ligament or augmentation with local tissues like fascia lata, often yielding poor results due to limited understanding of knee biomechanics and inadequate fixation techniques. The modern era of ACL reconstruction began in the 1960s when Dr. Kenneth Jones pioneered the use of the central third of the patellar tendon with bone blocks for ACL reconstruction, establishing the bone-patellar tendon-bone (BTB) autograft technique.

The introduction of arthroscopic surgery in the 1970s and 1980s, led by surgeons such as Dr. Richard O'Connor and Dr. Lanny Johnson, revolutionized ACL reconstruction by allowing for less invasive procedures, improved visualization, and reduced morbidity compared to traditional open techniques. The 1990s saw the popularization of hamstring tendon autografts, providing an alternative with potentially less anterior knee pain. Advances in fixation techniques transitioned from simple sutures to interference screws, cortical buttons, and cross-pins, enhancing graft incorporation and stability. The understanding of ACL anatomy progressed, leading to the development of "anatomic" single-bundle and later double-bundle reconstruction techniques, aiming to replicate the native ligament's biomechanics more closely. Ongoing research continues to refine surgical techniques, graft selection, and rehabilitation protocols, focusing on optimizing outcomes and facilitating a safe return to sport.

ACL reconstruction is indicated for patients presenting with symptomatic knee instability following an ACL tear, particularly those who wish to return to high-demand activities.

\begin{itemize}
  \item Complete ACL tear causing recurrent episodes of knee instability or "giving way" during daily activities or sports.
  \item ACL injury in young, active individuals (typically under 40-50 years old) who participate in cutting, pivoting, or jumping sports and desire to return to their pre-injury activity level.
  \item Combined ACL and meniscal tear (especially repairable meniscal tears) or other ligamentous injuries (e.g., medial collateral ligament, posterolateral corner), where instability from the ACL tear compromises healing or stability of the other structures.
  \item Persistent functional instability that significantly limits participation in desired recreational or occupational activities despite a trial of comprehensive non-operative rehabilitation.
  \item High-level athletes or individuals whose occupation demands a stable knee, where the risk of re-injury or further joint damage without reconstruction is unacceptable.
  \item Skeletally immature patients with an ACL tear and symptomatic instability, where specialized physeal-sparing techniques may be considered to prevent growth disturbance.
  \item Failure of conservative management, including bracing and physiotherapy, to adequately control symptoms of instability in an active individual.
\end{itemize}

ACL reconstruction is generally considered a primary, definitive surgical treatment for symptomatic ACL tears in active individuals. It is not an emergency procedure; rather, it is an elective surgery typically performed after an initial period of non-operative management focused on reducing swelling, restoring full range of motion, and regaining quadriceps strength. This pre-habilitation phase is crucial for optimizing postoperative outcomes and minimizing the risk of arthrofibrosis (knee stiffness).

For young, active patients who participate in pivoting sports, ACL reconstruction is often the first-line treatment recommendation due to the high likelihood of recurrent instability and subsequent meniscal or cartilage damage if left untreated. In older patients, or those with lower activity demands, a trial of conservative management with physical therapy and activity modification may be the initial approach. If symptomatic instability persists despite these measures, ACL reconstruction may then be considered as a secondary or salvage option. The decision to proceed with surgery is highly individualized, weighing the patient's age, activity level, associated injuries, and personal goals against the potential risks and benefits of the procedure.

\section*{Relevant Surgical Anatomy}
The knee joint is a complex synovial hinge joint formed by the articulation of the distal femur, proximal tibia, and patella. Key stabilizing structures include the menisci (medial and lateral), collateral ligaments (medial and lateral), and the cruciate ligaments (anterior and posterior). The anterior cruciate ligament (ACL) originates from the posteromedial aspect of the lateral femoral condyle, within the intercondylar notch, and inserts into the anterior intercondylar area of the tibia, just medial and anterior to the medial tibial spine. It comprises two functional bundles: the anteromedial (AM) bundle, which is taut in flexion, and the posterolateral (PL) bundle, which is taut in extension. The ACL's primary function is to prevent anterior translation of the tibia relative to the femur and to resist rotational forces.

Surrounding structures critical to ACL reconstruction include the infrapatellar fat pad, the synovium, and the neurovascular bundle located posteriorly in the popliteal fossa, which must be protected during tunnel drilling. Graft harvest requires specific anatomical considerations: for a patellar tendon autograft, the central third of the patellar tendon, along with bone blocks from the patella and tibia, is harvested. For hamstring autografts, the semitendinosus and gracilis tendons are harvested from their distal insertions on the pes anserinus, requiring careful identification and protection of the infrapatellar branch of the saphenous nerve and the saphenous vein. The quadriceps tendon autograft involves harvesting a portion of the quadriceps tendon, often with a patellar bone block. Accurate identification of the anatomical footprints of the native ACL on both the femur and tibia is paramount for successful graft placement and restoration of normal knee kinematics.

\section*{Preoperative Workup \& Preparation}
Thorough preoperative workup is essential for optimizing outcomes and minimizing complications. This typically begins with a comprehensive medical history and physical examination, focusing on knee stability, range of motion, and associated injuries. Imaging studies, primarily magnetic resonance imaging (MRI), are crucial to confirm the ACL tear, assess for meniscal or cartilage damage, and rule out other ligamentous injuries. Patients are encouraged to undergo a period of "pre-habilitation" to reduce swelling, regain full knee range of motion, and strengthen the quadriceps and hamstring muscles, as operating on a stiff or swollen knee significantly increases the risk of postoperative arthrofibrosis.

\begin{itemize}
  \item \textbf{Medical Clearance:} A general medical evaluation, including cardiac and pulmonary assessment, is performed to ensure the patient is fit for anesthesia and surgery.
  
  \item \textbf{Laboratory Tests:} Routine blood tests (e.g., complete blood count, coagulation profile, basic metabolic panel) are ordered.
  
  \item \textbf{Medication Review:} Anticoagulant and antiplatelet medications are typically discontinued or adjusted several days prior to surgery, as per anesthesiologist and surgeon recommendations.
  
  \item \textbf{NPO Status:} Patients are instructed to be NPO (nil per os) for a specified period (usually 6-8 hours) before surgery to prevent aspiration during anesthesia.
  
  \item \textbf{Prophylactic Antibiotics:} Intravenous prophylactic antibiotics are administered shortly before incision to reduce the risk of surgical site infection.
  
  \item \textbf{Informed Consent:} Detailed discussion with the patient regarding the procedure, graft options, potential risks, benefits, and expected postoperative course, including rehabilitation, is conducted to obtain informed consent.
  
  \item \textbf{Smoking Cessation:} Patients who smoke are strongly advised to cease smoking several weeks before surgery, as it impairs wound healing and graft incorporation.
\end{itemize}

Careful patient selection is paramount for successful ACL reconstruction.

\textbf{Absolute Contraindications:}
\begin{itemize}
  \item Active infection within the knee joint or surrounding tissues.
  
  \item Severe, uncontrolled systemic medical conditions that preclude safe anesthesia or surgery (e.g., unstable angina, recent myocardial infarction, severe coagulopathy).
  
  \item Inability or unwillingness of the patient to participate in a rigorous postoperative rehabilitation program.
\end{itemize}

\textbf{Relative Contraindications \& Precautions:}
\begin{itemize}
  \item \textbf{Advanced Osteoarthritis:} Significant pre-existing knee osteoarthritis may lead to suboptimal functional outcomes, as the underlying degenerative process will continue. In such cases, the benefits of reconstruction must be carefully weighed against the risk of persistent pain and limited improvement.
  
  \item \textbf{Severe Knee Stiffness or Flexion Contracture:} Operating on a knee with limited range of motion or a significant flexion contracture increases the risk of postoperative arthrofibrosis. Preoperative rehabilitation to restore motion is crucial.
  
  \item \textbf{Advanced Age with Low Activity Demands:} In older, less active individuals, conservative management with physical therapy and bracing may be more appropriate if symptoms are manageable, as the functional benefit of surgery may not outweigh the risks.
  
  \item \textbf{Skeletally Immature Patients:} In children and adolescents with open growth plates, traditional tunnel drilling techniques carry a risk of physeal injury and growth disturbance. Physeal-sparing or transphyseal techniques with careful consideration are necessary.
  
  \item \textbf{Obesity:} While not an absolute contraindication, severe obesity can increase surgical complexity, anesthetic risks, and potentially impact rehabilitation and graft healing.
  
  \item \textbf{Poor Bone Quality:} Osteoporosis or other conditions leading to poor bone quality can compromise graft fixation.
  
  \item \textbf{Revision Surgery:} Previous failed ACL reconstruction presents unique challenges, often requiring more complex techniques, different graft choices, and careful assessment of tunnel placement.
\end{itemize}

\section*{Surgical Technique \& Procedure}
ACL reconstruction is typically performed under general or regional anesthesia (e.g., spinal or epidural) and is often an outpatient procedure. The patient is positioned supine on the operating table, and a tourniquet is applied to the proximal thigh to minimize bleeding.

\begin{itemize}
  \item \textbf{Arthroscopic Evaluation:} Small portal incisions are made around the knee to insert an arthroscope and surgical instruments. The knee joint is thoroughly inspected to confirm the ACL tear, assess the condition of the menisci and articular cartilage, and identify any associated injuries. Meniscal tears are addressed at this stage, either by repair or partial meniscectomy.
  
  \item \textbf{Graft Harvest:} The chosen autograft (e.g., central third of the patellar tendon with bone blocks, semitendinosus and gracilis tendons, or quadriceps tendon with or without a bone block) is harvested through a separate, small incision. If an allograft is used, it is prepared from sterile donor tissue. The graft is then prepared to the appropriate length and diameter, often folded or braided, and sutures are placed at its ends for handling and fixation.
  
  \item \textbf{Tunnel Preparation:} The remnants of the torn ACL are debrided. Under arthroscopic guidance, bone tunnels are precisely drilled in the femur and tibia at the anatomical attachment sites of the native ACL. The femoral tunnel is typically created through an accessory medial portal or a transtibial approach, while the tibial tunnel is drilled from the anteromedial tibia. Careful attention is paid to tunnel position and angulation to ensure isometric graft placement.
  
  \item \textbf{Graft Passage:} A guide wire is passed through the tibial tunnel, across the joint, and out through the femoral tunnel. The prepared graft is then pulled through the tibial and femoral tunnels using a strong suture, ensuring it is properly seated within the joint.
  
  \item \textbf{Graft Fixation:} Once the graft is in the correct position and under appropriate tension (often with the knee in a specific degree of flexion), it is secured in both the femoral and tibial tunnels using various fixation devices. Common methods include interference screws (bioabsorbable or metallic), cortical buttons (e.g., Endobutton), cross-pins, or staples. The fixation aims to provide strong initial stability to allow for biological integration.
  
  \item \textbf{Final Assessment \& Closure:} The knee's range of motion and stability are assessed to confirm proper graft tension and absence of impingement. The arthroscopic portals and graft harvest incisions are closed in layers, typically with absorbable sutures for deeper layers and non-absorbable sutures or staples for the skin. A sterile dressing is applied, and the knee is often placed in a brace.
\end{itemize}

\section*{Complications \& Risks}
As with any surgical procedure, ACL reconstruction carries potential risks and complications, which can be broadly categorized as general surgical risks and procedure-specific risks.

\textbf{General Surgical Complications:}
\begin{itemize}
  \item \textbf{Anesthesia-related risks:} Adverse reactions to anesthetic agents, respiratory complications, cardiovascular events.
  
  \item \textbf{Infection:} Surgical site infection (superficial or deep intra-articular infection), which can lead to graft failure and require further surgery.
  
  \item \textbf{Bleeding/Hematoma:} Excessive bleeding during or after surgery, potentially leading to hemarthrosis (blood in the joint) requiring aspiration.
  
  \item \textbf{Deep Vein Thrombosis (DVT) / Pulmonary Embolism (PE):} Blood clot formation in the leg veins, with the risk of migration to the lungs.
  
  \item \textbf{Nerve Injury:} Damage to peripheral nerves (e.g., saphenous nerve, infrapatellar branch of the saphenous nerve, common peroneal nerve) leading to numbness, weakness, or pain.
  
  \item \textbf{Vascular Injury:} Damage to major blood vessels, though rare, can lead to significant bleeding or limb ischemia.
\end{itemize}

\textbf{Procedure-Specific Complications:}
\begin{itemize}
  \item \textbf{Graft Failure or Re-tear:} The most significant risk, where the reconstructed ligament stretches, loosens, or tears again, often requiring revision surgery. Risk factors include young age, return to high-risk sports, and technical errors.
  
  \item \textbf{Persistent Knee Instability:} Despite an intact graft, some patients may experience residual laxity or instability if tunnel placement is suboptimal or if there are other untreated ligamentous injuries.
  
  \item \textbf{Arthrofibrosis (Knee Stiffness):} Formation of excessive scar tissue within the joint, leading to loss of knee extension or flexion, often requiring intensive physical therapy or manipulation under anesthesia.
  
  \item \textbf{Anterior Knee Pain:} More common after patellar tendon autograft harvest, due to patellar tendonitis or patellofemoral pain.
  
  \item \textbf{Donor Site Morbidity:} Pain, weakness, or numbness at the site where the autograft was harvested (e.g., hamstring weakness, quadriceps weakness, infrapatellar numbness).
  
  \item \textbf{Hardware Complications:} Irritation from fixation devices, screw breakage, or migration, potentially requiring removal.
  
  \item \textbf{Cyclops Lesion:} A nodule of fibrous tissue that can form in the intercondylar notch, causing mechanical impingement and loss of extension.
  
  \item \textbf{Long-term Osteoarthritis:} While reconstruction aims to reduce this risk, ACL injury itself is a significant risk factor for developing knee osteoarthritis, regardless of surgical intervention.
  
  \item \textbf{Growth Plate Injury:} In skeletally immature patients, improper tunnel placement can damage the physis, leading to growth arrest or deformity.
\end{itemize}

\section*{Postoperative Care \& Recovery}
The postoperative recovery from ACL reconstruction is a structured, multi-phase process that typically spans 6 to 12 months, with full return to sport often taking up to a year or more.

Immediately after surgery, the patient is monitored in the Post-Anesthesia Care Unit (PACU) for recovery from anesthesia, pain control, and neurovascular status of the limb. The knee is typically protected in a brace, often locked in extension, and crutches are used for initial ambulation. The first 24-48 hours focus on pain management, controlling swelling with ice and elevation, and initiating gentle range-of-motion exercises, particularly emphasizing full knee extension to prevent stiffness. Weight-bearing status is determined by the surgeon and graft type, often partial weight-bearing initially, progressing to full weight-bearing over the first few weeks.

Over the first 1-2 weeks, the primary goals are to achieve full knee extension, minimize swelling, and begin quadriceps activation exercises. Physical therapy sessions become more frequent, focusing on restoring normal gait mechanics, improving flexion, and initiating light strengthening. Diet is advanced as tolerated, and activity restrictions include avoiding twisting or pivoting motions. Long-term recovery, spanning several months, involves progressive strengthening of the quadriceps, hamstrings, and gluteal muscles, proprioceptive training, and sport-specific drills. Return to running is typically permitted around 3-4 months, and return to cutting and pivoting sports is usually not before 6-9 months, contingent upon meeting specific functional criteria, strength benchmarks, and surgeon clearance.

During ACL reconstruction, the surgeon meticulously documents several key findings. These include the precise location and extent of the ACL tear (e.g., proximal, mid-substance, distal avulsion), the quality of the remaining ACL stump, and the condition of the intercondylar notch. Crucially, any associated injuries to the menisci (medial or lateral), articular cartilage, or other ligaments (e.g., MCL, LCL) are identified and addressed. The type and severity of meniscal tears, if present, are noted, along with whether they were repaired or debrided. The presence of loose bodies or signs of early osteoarthritis are also recorded.

Pathology reports on resected tissues, such as meniscal fragments or ACL remnants, confirm the histological nature of the tissue. For instance, meniscal biopsies confirm fibrocartilaginous tissue, and ACL remnants confirm ligamentous tissue. While the primary graft itself is not typically sent for pathological analysis unless there is a specific concern (e.g., infection), the pathology report provides definitive confirmation of any associated injuries that were treated, contributing to a comprehensive understanding of the patient's knee pathology.

A normal and successful postoperative course following ACL reconstruction is characterized by a progressive and predictable return of knee function and stability. Initially, patients can expect some pain and swelling, which should gradually subside with medication, ice, and elevation. Within the first few weeks, pain should become manageable with over-the-counter analgesics, and the surgical incisions should heal without complications. The knee's range of motion should steadily improve, with the goal of achieving full extension and near-full flexion within the first 6-8 weeks.

Patients should experience a gradual restoration of quadriceps and hamstring strength, leading to improved gait and the ability to perform daily activities without crutches. There should be no sensation of "giving way" or instability. As rehabilitation progresses, patients will regain confidence in their knee, with improved balance, proprioception, and the ability to perform more demanding physical tasks. A successful course culminates in a stable, functional knee that allows the patient to return to their desired activity level, including sports, without pain or instability, typically within 9-12 months, provided they adhere strictly to the rehabilitation protocol.

Postoperative care for ACL reconstruction is centered around a structured and progressive rehabilitation program, guided by regular follow-up with the orthopedic surgeon and physical therapist.

\begin{itemize}
  \item \textbf{Scheduled Postoperative Visits:} Patients typically have follow-up appointments at 1-2 weeks, 6 weeks, 3 months, 6 months, and 1 year post-surgery to monitor incision healing, assess knee range of motion and stability, and evaluate progress in rehabilitation.
  
  \item \textbf{Physical Therapy (PT):} A comprehensive, individualized physical therapy program is initiated immediately after surgery and is crucial for optimal recovery. It progresses through phases focusing on pain and swelling control, range of motion, quadriceps activation, strengthening, proprioception, and sport-specific training.
  
  \item \textbf{Suture/Staple Removal:} Skin sutures or staples are usually removed at the first postoperative visit, typically 10-14 days after surgery.
  
  \item \textbf{Brace and Crutch Use:} The use of a knee brace and crutches is prescribed based on surgeon preference and graft type, with gradual weaning as stability and strength improve.
  
  \item \textbf{Functional Testing:} Before returning to high-impact activities or sports, patients undergo objective functional testing (e.g., hop tests, strength assessments) to ensure adequate strength, stability, and neuromuscular control.
  
  \item \textbf{Warning Signs:} Patients are educated on warning signs that necessitate immediate contact with their surgeon, including fever, chills, persistent or worsening pain, redness or drainage from incisions, calf pain or swelling, or sudden loss of knee motion.
  
  \item \textbf{Long-term Monitoring:} Patients are advised about the long-term risk of osteoarthritis and the importance of maintaining knee health through exercise and activity modification.
\end{itemize}

\section*{Outcomes \& Prognosis}
Anterior cruciate ligament injuries are among the most prevalent and debilitating knee injuries, particularly affecting athletic populations. The estimated incidence in the United States ranges from 100,000 to 200,000 cases per year, with ACL reconstruction being one of the most frequently performed orthopedic procedures. These injuries occur most commonly in young, active individuals, with a peak incidence in those aged 15-25 years. Females are disproportionately affected, exhibiting a 2-8 times higher incidence rate than males in similar cutting, pivoting, and jumping sports such as soccer, basketball, volleyball, and skiing. The majority of ACL tears (approximately 70-80\%) are non-contact injuries, typically occurring during sudden deceleration, landing from a jump, or rapid change of direction. While mortality directly associated with ACL injury is negligible, the morbidity is significant, encompassing chronic knee instability, impaired athletic performance, and a substantially increased risk of secondary meniscal tears and articular cartilage damage if left untreated.

ACL reconstruction generally yields favorable outcomes, with a high success rate in restoring knee stability and enabling a return to pre-injury activity levels for the majority of patients. Reported success rates, defined by a stable knee and satisfactory functional outcome, range from 85\% to 95\%. Return to sport rates are also high, with 60-80\% of athletes returning to their pre-injury level of competition, although this can vary based on sport and patient expectations. Graft failure or re-tear rates are typically low, often below 10\% in the general population, but can be higher (up to 20-30\%) in very young, highly active athletes who return to high-risk sports.

The long-term prognosis, however, is complicated by the inherent risk of post-traumatic osteoarthritis (PTOA) following an ACL injury, regardless of surgical intervention. Studies have shown that even with successful reconstruction, a significant proportion of patients may develop radiographic signs of osteoarthritis within 10-15 years. Prognostic factors influencing outcomes include patient age, activity level, presence of associated meniscal or cartilage injuries, graft choice (e.g., autograft generally has lower re-tear rates than allograft in young athletes), surgical technique (e.g., anatomical tunnel placement), and strict adherence to a comprehensive rehabilitation program. Landmark studies, such as those comparing different graft types or evaluating long-term outcomes, consistently highlight the importance of restoring stability to protect the knee, while also acknowledging the persistent challenge of PTOA.

\section*{References}
\begin{enumerate}
  \item American Academy of Orthopaedic Surgeons. ACL Injury: Diagnosis and Treatment. AAOS; 2024.
  
  \item Frank RM, Van Eck CF, Brand J, et al. Anatomic ACL Reconstruction: A Cadaveric Study Comparing Double-Bundle and Single-Bundle Techniques. Arthroscopy. 2012;28(10):1479-1487.
  
  \item Marx RG, Jones EC, Allen AA, et al. The effect of graft choice on the outcome of anterior cruciate ligament reconstruction. Am J Sports Med. 2005;33(10):1491-1498.
  
  \item Sabiston Textbook of Surgery. 22nd ed. Elsevier; 2022.
\end{enumerate}

\clearpage
